%-------------------------
% CV - Johan Carlo Amador Egoavil
% Overleaf / LaTeX
%-------------------------

\documentclass[a4paper,10pt]{article}

\usepackage[utf8]{inputenc}
\usepackage[T1]{fontenc}
\usepackage[margin=0.65in]{geometry}
\usepackage{titlesec}
\usepackage{enumitem}
\usepackage{hyperref}
\usepackage{xcolor}
\usepackage{tabularx}
\usepackage{fontawesome5}
\setlength{\parindent}{0pt}
\setlength{\emergencystretch}{2em}

% Colors
\definecolor{primary}{HTML}{171717}
\definecolor{darktext}{HTML}{1a1a1a}
\definecolor{graytext}{HTML}{555555}

% Hyperlinks
\hypersetup{
    colorlinks=true,
    hypertexnames=false,
    linkcolor=primary,
    urlcolor=primary,
}

% Section formatting
\titleformat{\section}{\large\bfseries\color{darktext}}{}{0em}{}[\titlerule]
\titlespacing*{\section}{0pt}{9pt}{4pt}

% Remove page numbers
\pagenumbering{gobble}

% Custom commands
\newcommand{\cvsubheading}[4]{
  \vspace{2pt}
  \begin{tabularx}{\textwidth}{@{}X r@{}}
    \textbf{#1} & \textit{\small #2} \\
    \textit{\small\color{graytext} #3} & \textit{\small\color{graytext} #4} \\
  \end{tabularx}\vspace{-2pt}
}

\newcommand{\cvproject}[3]{
  \vspace{2pt}
  \begin{tabularx}{\textwidth}{@{}X r@{}}
    \textbf{#1} & \textit{\small #2} \\
  \end{tabularx}
  {\small\color{graytext} #3}\vspace{-2pt}
}

\begin{document}

%----------HEADER----------
\begin{center}
  {\LARGE\bfseries Johan Carlo Amador Egoavil} \\[3pt]
  {\small\color{graytext} Desarrollo Full-Stack | Interoperabilidad en Salud} \\[6pt]
  \small
  \faIcon{map-marker-alt}\enspace Lima, Per\'u \quad
  \faIcon{envelope}\enspace \href{mailto:johan.amador@pucp.edu.pe}{johan.amador@pucp.edu.pe} \quad
  \faIcon{phone}\enspace +51 951 665 323 \\[3pt]
  \faIcon{globe}\enspace \href{https://johanamador.com}{johanamador.com} \quad
  \faIcon{github}\enspace \href{https://github.com/johanamador}{johanamador} \quad
  \faIcon{linkedin}\enspace \href{https://www.linkedin.com/in/johanamadordev/}{johanamadordev}
\end{center}

%----------PROFILE----------
\section{Perfil Profesional}
Egresado de Ingeniería Informática de la PUCP con experiencia en desarrollo full-stack, sistemas hospitalarios e interoperabilidad en salud. En PUCP / GIDIS desarrollo los flujos de admisión de SIH.SALUS y un módulo de interoperabilidad HL7 FHIR / IPS para RENHICE. Desarrollo plataformas web, APIs y sistemas internos desde el análisis de requisitos hasta el despliegue, capacitación y mantenimiento, incluso en entornos con conectividad limitada.

%----------EXPERIENCE----------
\section{Experiencia Profesional}

\cvsubheading
  {Desarrollador de Software -- Sistemas de Salud}{Agosto 2025 -- Actualidad}
  {PUCP / GIDIS -- SIH.SALUS}{H\'ibrido, Lima}
\begin{itemize}[leftmargin=16pt, itemsep=1.5pt, parsep=0pt, topsep=3pt]
  \item Desarrollo e implementaci\'on del sistema de historias cl\'inicas electr\'onicas del Hospital Santa Clotilde, como parte del Grupo de Investigaci\'on y Desarrollo en Ingenier\'ia de Software (GIDIS). Responsable del m\'odulo de admisi\'on e integraci\'on de est\'andares de interoperabilidad en salud.
  \item Despliegue presencial en agosto de 2026 en Santa Clotilde, cuenca del r\'io Napo: admisi\'on, triaje, farmacia, laboratorio y consulta externa. Traslado e instalaci\'on del servidor y una UPS con el equipo, en un hospital con energ\'ia y conectividad inestables.
  \item Capacitaci\'on al personal de admisi\'on, seguimiento del uso del sistema y mejoras posteriores al despliegue.
\end{itemize}


\cvsubheading
  {Desarrollador Frontend}{Junio 2026 -- Actualidad}
  {NexoDev -- Contrato}{Remoto}
\begin{itemize}[leftmargin=16pt, itemsep=1.5pt, parsep=0pt, topsep=3pt]
  \item Desarrollo de interfaces web en producción e integración con APIs y servicios backend existentes.
  \item Colaboración remota en nuevas funcionalidades, resolución de errores y mejoras iterativas del producto.
\end{itemize}

\cvsubheading
  {Ingeniero de Software (Freelance)}{Noviembre 2024 -- Actualidad}
  {Desarrollo full stack -- Gesti\'on integral de proyectos}{Lima, Per\'u}

\begin{itemize}[leftmargin=16pt, itemsep=1.5pt, parsep=0pt, topsep=3pt]
  \item Gesti\'on aut\'onoma end-to-end: an\'alisis de requisitos, dise\~no t\'ecnico, desarrollo full-stack, despliegue en VPS/hosting del cliente, capacitaci\'on y mantenimiento continuo.
  \item Desarrollo de frontends en Next.js y backends/API en PHP + MySQL para cat\'alogos, paneles administrativos, carga de im\'agenes, autenticaci\'on simple y flujos operativos.
  \item Ciclos iterativos de feedback, correcci\'on de errores, optimizaci\'on de rendimiento y mejoras funcionales post-lanzamiento.
\end{itemize}

\vspace{2pt}
{\small\textbf{Proyectos destacados para clientes:}}
\begin{itemize}[leftmargin=16pt, itemsep=1.5pt, parsep=0pt, topsep=2pt]
  \item \textbf{\href{https://gruposercom.pe}{Grupo Sercom} / \href{https://intranet.gruposercom.pe}{Intranet}} -- Web B2B y portal documental con acceso por RUC, vista previa PDF y administración privada. Next.js, TypeScript, VPS, PM2.
  \item \textbf{\href{https://prosedain.com}{Prosedain}} -- Catálogo industrial, administración de productos y cotizaciones PDF. Next.js, PHP, MySQL.
  \item \textbf{\href{https://www.isermaperu.com}{Iserma}} -- Catálogo industrial con productos, marcas y categorías editables. Next.js, React, Mantine, PHP, MySQL.
  \item \textbf{\href{https://mecanicalevano.org}{ERP Automotriz Lévano}} -- Plataforma privada para la empresa automotriz con acceso autenticado del equipo. TypeScript, React, Next.js.
  \item \textbf{\href{https://flow-telligence.com}{Nebu}} -- Voz y video en tiempo real, integraciones de IA y pagos en línea. Remix, NestJS, PostgreSQL, Docker, LiveKit.
\end{itemize}
{\small\color{graytext} Otros proyectos: \href{https://farmasaludinversiones.com}{Farmasalud Inversiones}, \href{https://mitsperu.com}{MITS Perú}, \href{https://academiapasalo.com}{Academia Pásalo}.}

\section{Interoperabilidad en Salud}
\textbf{Conectatón IPS Perú 2025} -- Participación técnica en representación del Hospital Santa Clotilde y la PUCP en intercambio de información clínica con estándares internacionales.

\newpage
%----------THESIS----------
\section{Tesis}

\cvsubheading
  {SIH.SALUS -- Módulo de Interoperabilidad en Salud}{}
  {Pontificia Universidad Cat\'olica del Per\'u}{2025 -- 2026}

\begin{itemize}[leftmargin=16pt, itemsep=1.5pt, parsep=0pt, topsep=3pt]
  \item M\'odulo de interoperabilidad para OpenMRS integrado con RENHICE (Registro Nacional de Historias Cl\'inicas Electr\'onicas del MINSA) usando HL7 FHIR.
  \item Mapeo de datos clínicos de OpenMRS a recursos FHIR y bundles IPS, con envío de resúmenes y consulta/importación bidireccional.
  \item Colas persistentes, procesamiento diferido, reintentos y monitoreo para conectividad intermitente; pruebas de integración con HAPI FHIR.
\end{itemize}
{\scriptsize\color{graytext} OpenMRS 3 $\cdot$ Java $\cdot$ HL7 FHIR $\cdot$ Interoperabilidad cl\'inica} \quad {\small \faIcon{github}\enspace \href{https://github.com/sihsalus/openmrs-module-sihsalusinterop}{Mi módulo de interoperabilidad}}

%----------ACADEMIC PROJECTS----------
\section{Proyectos Universitarios}

\cvproject{Sistema de Gesti\'on de Tesis}{Marzo 2025 -- Julio 2025}{React $\cdot$ Spring Boot $\cdot$ PostgreSQL}
\par\smallskip Proyecto final con 38 integrantes. Desarrollo del módulo de revisión e integración de GoWinston para detección de plagio; participación en CI/CD, control de calidad y revisión de código.\par\medskip

\cvproject{Sistema de Planificaci\'on de Operaciones Log\'isticas}{Agosto 2024 -- Diciembre 2024}{Java $\cdot$ React}
\par\smallskip Planificación de rutas con restricciones y bloqueos, motor de optimización en Java e interfaz modular en React.\par\medskip

\cvproject{Sistema de Gesti\'on de Tutor\'ias Acad\'emicas}{Marzo 2024 -- Julio 2024}{React $\cdot$ ASP.NET}
\par\smallskip Plataforma de tutorías universitarias con horarios, notificaciones y gestión de sesiones, desarrollada con Scrum.\par\medskip

\cvproject{Sistema de Gesti\'on de Riesgos Bancarios}{Agosto 2023 -- Diciembre 2023}{Java $\cdot$ C\#}
\par\smallskip Prototipo académico de riesgos bancarios con registro, análisis, jerarquía y reportes con filtros.\par\medskip

%----------EDUCATION----------
\section{Educaci\'on}

\cvsubheading
  {Ingenier\'ia Inform\'atica (Egresado)}{2020 -- 2026}
  {Pontificia Universidad Cat\'olica del Per\'u (PUCP)}{Lima, Per\'u}

\vspace{2pt}
\cvsubheading
  {Full Stack Developer Path}{2022 -- 2023}
  {Platzi}{Online}

\vspace{2pt}
\cvsubheading
  {Software Development}{2025}
  {Coursera}{Online}

%----------SKILLS----------
\section{Habilidades T\'ecnicas}

\begin{tabularx}{\textwidth}{@{}l X@{}}
  \textbf{Lenguajes:} & Java, JavaScript, TypeScript, PHP, C/C++, C\#, Python, SQL \\[3pt]
  \textbf{Frameworks:} & React, Next.js, NestJS, Spring Boot, ASP.NET, Express.js, Remix, Mantine \\[3pt]
  \textbf{Bases de Datos:} & PostgreSQL, MySQL, SQL Server, Oracle SQL, MongoDB, Supabase \\[3pt]
  \textbf{DevOps \& Cloud:} & Docker, CI/CD, VPS, PM2, Vercel, cPanel, AWS (RDS), Namecheap, Cloudflare, Linux, Git/GitHub \\[3pt]
  \textbf{Herramientas:} & Postman, Jira, VS Code, Notion, Bizagi (BPMN) \\[3pt]
  \textbf{Salud Digital:} & HL7 FHIR, HL7 v2, IPS, OpenMRS, HAPI FHIR, mapeo de datos clínicos \\[3pt]
  \textbf{APIs e Integración:} & REST APIs, GraphQL, Prisma, NextAuth, Cloudinary, generación PDF \\[3pt]
  \textbf{Interfaces:} & Tailwind CSS, shadcn/ui \\[3pt]
  \textbf{Metodolog\'ias:} & Scrum, desarrollo \'agil, auditor\'ia t\'ecnica, despliegue continuo, control de versiones, trabajo colaborativo, documentaci\'on t\'ecnica \\
\end{tabularx}

%----------LANGUAGES----------
\section{Idiomas}

\begin{tabular}{@{}l l}
  \textbf{Espa\~nol:} & Nativo \\[3pt]
  \textbf{Ingl\'es:} & Intermedio-Avanzado (B2 -- FCE Cambridge) \\
\end{tabular}

\end{document}
